% Generated draft: requires mathematical and source review.
\documentclass[11pt]{article}
\usepackage[margin=25mm]{geometry}
\usepackage{fontspec}
\setmainfont{Latin Modern Roman}
\usepackage{amsmath,amssymb,mathtools,unicode-math}
\setmathfont{Latin Modern Math}
\usepackage{enumitem,needspace,xurl,hyperref}
\setlength{\parindent}{0pt}
\setlength{\parskip}{5pt}
\setlength{\emergencystretch}{3em}
\newcommand{\N}{\mathbb{N}}
\newcommand{\Z}{\mathbb{Z}}
\newcommand{\Q}{\mathbb{Q}}
\newcommand{\R}{\mathbb{R}}
\newcommand{\C}{\mathbb{C}}
\newcommand{\F}{\mathbb{F}}
\newcommand{\A}{\mathbb{A}}
\newcommand{\PP}{\mathbb{P}}
\newcommand{\E}{\mathbb{E}}
\newcommand{\eps}{\varepsilon}
\newcommand{\dd}{\,\mathrm{d}}
\begin{document}
\section*{UnsolvedMath ID: 701006}
\textbf{Source problem code:} catalogue entry
\subsection{Definitions and mathematical statement}
Let \(p\) be a prime. A pro-\(p\) group is a topological group isomorphic to an inverse limit of finite \(p\)-groups. Let \(F\) be a free pro-\(p\) group, with freeness understood in the category of pro-\(p\) groups and continuous homomorphisms.

The supplied record does not define the Noetherian Equation Property. Under the standard interpretation as equational Noetherianity with coefficients, its meaning is as follows. For each integer \(n\geq 1\), let \(L_n\) be the abstract free group on \(x_1,\ldots,x_n\). An equation over \(F\) is an expression \(w(x_1,\ldots,x_n)=1\), where \(w\in F*L_n\) and \(*\) denotes the abstract free product. For a system \(S\subseteq F*L_n\), define
\[
V_F(S)=\{(a_1,\ldots,a_n)\in F^n: w(a_1,\ldots,a_n)=1\text{ for every }w\in S\}.
\]
Here evaluation fixes the coefficients in \(F\). The property is
\[
\forall n\geq 1\quad\forall S\subseteq F*L_n\quad
\exists S_0\subseteq S\quad
\bigl(|S_0|<\infty\ \text{ and }\ V_F(S_0)=V_F(S)\bigr).
\]
The question is whether a free pro-\(p\) group has this property.

This is a conditional rendering of the source terminology: the supplied record does not specify whether equations permit coefficients, whether only ordinary finite group words are intended, or whether a notion involving pro-\(p\) completions is intended. It also does not specify a rank restriction or explicitly state the quantification over primes and ranks. Consequently, the displayed property gives a precise conventional interpretation, while the intended scope of the original question remains unresolved.
\subsection{Short English statement}
Under the conventional interpretation, must every system of equations with coefficients in a free pro-\(p\) group, in finitely many variables, have the same solutions as some finite subsystem? The source leaves the equation convention and the intended range of primes and ranks unspecified.
\subsection{Sources}
\begin{itemize}
\item Supplied local UnsolvedMath registry record, permanent ID \(701006\), exact source problem code \texttt{catalogue entry}, titled ``Kourovka Notebook Problem 14.21''; upstream revision \texttt{d371ee917ac9647fbb06cd6fc529deec6f3dda82}.
\item The supplied attribution identifies G. Baumslag, A. G. Myasnikov, and V. N. Remeslennikov as the proposers, with proposed year 1999.
\item The supplied background cites \emph{The Kourovka Notebook (Unsolved Problems in Group Theory)}, 14th issue (1999), Problem 14.21, and the 21st issue edition (2026), \texttt{arXiv:1401.0300v48}. The collection description names E. I. Khukhro and V. D. Mazurov as editors. These bibliographic details are supplied attribution and have not been independently checked.
\item The registry labels the problem ``open''. This is a supplied status label, not an independently verified mathematical or literature assessment.
\end{itemize}
\end{document}
